\documentclass[../../main.tex]{subfiles}
\graphicspath{{../../image/}{../../figure/}}
\begin{document}

\section{扩散效应}\label{section:ch15-扩散效应}

正如\cref{section:ch15-证明路线图}中所指出，黏性的作用是在某个长度为 $\bar t=O(\delta_0^{-2})$ 的时间区间 $(0,\bar t]$ 上维持\eqref{align:ch15-1523}，同时如\eqref{align:ch15-1524}所示地减小高阶空间导数的 $L^1$ 范数。这两项目标均由下述引理实现。

\begin{lemma}\label{lemma:ch15-扩散效应估计}
设 $U$ 是\eqref{align:ch15-1521}、\eqref{align:ch15-1522}的解。存在正常数 $\alpha$ 与 $\kappa$，使得对任意固定的充分小正数 $\delta_0$，若\eqref{align:ch15-1523}在长度为 $\bar t=(\alpha\kappa\delta_0)^{-2}$ 的区间 $(0,\bar t]$ 上成立，则
\begin{align}
\|\partial_x^2U(\cdot,t)\|_{L^1(-\infty,\infty)}&<\frac{2\delta_0}{\kappa\sqrt{t}},\qquad t\in(0,\bar t],
\label{align:ch15-1531}\\
\|\partial_x^3U(\cdot,t)\|_{L^1(-\infty,\infty)}&<\frac{5\delta_0}{\kappa^2t},\qquad t\in(0,\bar t],
\label{align:ch15-1532}\\
\|\partial_x^3U(\cdot,t)\|_{L^\infty(-\infty,\infty)}&<\frac{16\delta_0}{\kappa^3t\sqrt{t}},\qquad t\in(0,\bar t].
\label{align:ch15-1533}
\end{align}
进而，当\eqref{align:ch15-1523}在更长的区间 $(0,T)$ 上成立、$\bar t<T\leqslant\infty$ 时，\eqref{align:ch15-1524}对任意 $t\in[\bar t,T)$ 成立。最后，
\begin{align}
TV_{(-\infty,\infty)}U_0(\cdot)<\frac12\kappa\delta_0
\label{align:ch15-1534}
\end{align}
蕴含\eqref{align:ch15-1523}对一切 $t\in(0,\bar t]$ 成立。
\end{lemma}

\begin{proof}[证明梗概]
把\eqref{align:ch15-1521}改写为
\begin{align}
\partial_tU+A^*\partial_xU-\partial_x^2U=[A^*-A(U)]\partial_xU.
\label{align:ch15-1535}
\end{align}
\eqref{align:ch15-1535}左端线性抛物算子的 $(n\times n)$ 矩阵值 Green 核 $G(x,t)$ 可写成闭形式。将\eqref{align:ch15-1535}左端乘以 $L_i^*$，便把该系统解耦为标量方程，其算子为 $\partial_t+\lambda_i^*\partial_x-\partial_x^2$，其 Green 函数为
\begin{align}
g_i(x,t)=\frac1{2\sqrt{\pi t}}\exp\left[-\frac{(x-\lambda_i^*t)^2}{4t}\right].
\label{align:ch15-1536}
\end{align}
因此，
\begin{align}
G(x,t)=\sum_{i=1}^ng_i(x,t)R_i^*L_i^*.
\label{align:ch15-1537}
\end{align}
简单计算给出
\begin{align}
\|G(\cdot,t)\|_{L^1(-\infty,\infty)}\leqslant\frac1\kappa,\qquad\|\partial_xG(\cdot,t)\|_{L^1(-\infty,\infty)}\leqslant\frac1{\kappa\sqrt{t}},
\label{align:ch15-1538}
\end{align}
其中 $\kappa$ 为某常数。

在附加假设 $U_0\in C^\infty$ 下建立所需估计即可，因为一般情形随后通过完备化得到。

对 $x$ 微分\eqref{align:ch15-1535}，并对所得方程应用 Duhamel 原理，得到
\begin{align}
\partial_xU(\cdot,t)=G(\cdot,t)*\partial_xU_0(\cdot)+\int_0^tG(\cdot,t-\tau)*P(\cdot,\tau)\,\mathrm{d}\tau,
\label{align:ch15-1539}
\end{align}
其中
\begin{align}
P(x,\tau)=[A^*-A(U(x,\tau))]\partial_x^2U(x,t)-\partial_xU^\top(x,\tau)DA(U(x,\tau))\partial_xU(x,\tau)
\label{align:ch15-15310}
\end{align}
而 $*$ 表示关于变量 $x$ 在 $(-\infty,\infty)$ 上的卷积。

由 $\|U-U^*\|_{L^\infty}\leqslant\|\partial_xU\|_{L^1}$ 与 $\|\partial_xU\|_{L^\infty}\leqslant\|\partial_x^2U\|_{L^1}$，
\begin{align}
\|P(\cdot,\tau)\|_{L^1(-\infty,\infty)}\leqslant\beta\|\partial_xU(\cdot,\tau)\|_{L^1(-\infty,\infty)}\|\partial_x^2U(\cdot,\tau)\|_{L^1(-\infty,\infty)},
\label{align:ch15-15311}
\end{align}
其中 $\beta$ 仅依赖于 $\sup_{U\in\mathscr{O}}|D^2F(U)|$。

现假设\eqref{align:ch15-1523}在长度为 $\bar t=(\alpha\kappa\delta_0)^{-2}$ 的区间 $[0,\bar t]$ 上成立，$\alpha>2\pi\beta\kappa^{-2}$，并在 $(0,\bar t]$ 上验证\eqref{align:ch15-1531}。对 $x$ 微分\eqref{align:ch15-1539}，
\begin{align}
\partial_x^2U(\cdot,t)=\partial_xG(\cdot,t)*\partial_xU_0(\cdot)+\int_0^t\partial_xG(\cdot,t-\tau)*P(\cdot,\tau)\,\mathrm{d}\tau.
\label{align:ch15-15312}
\end{align}
联合\eqref{align:ch15-1538}、\eqref{align:ch15-15311}与\eqref{align:ch15-1523}，这蕴含
\begin{align}
\|\partial_x^2U(\cdot,t)\|_{L^1(-\infty,\infty)}\leqslant\frac{\delta_0}{\kappa\sqrt{t}}+\frac{\beta\delta_0}{\kappa}\int_0^t\frac1{\sqrt{t-\tau}}\|\partial_x^2U(\cdot,\tau)\|_{L^1(-\infty,\infty)}\,\mathrm{d}\tau.
\label{align:ch15-15313}
\end{align}

设\eqref{align:ch15-1531}不成立，并令 $t$ 为 $(0,\bar t]$ 中它首次失效的时刻。则由\eqref{align:ch15-15313}得
\begin{align}
\frac{2\delta_0}{\kappa\sqrt{t}}\leqslant\frac{\delta_0}{\kappa\sqrt{t}}+\frac{2\beta\delta_0^2}{\kappa^2}\int_0^t\frac1{\sqrt{\tau(t-\tau)}}\,\mathrm{d}\tau=\frac{\delta_0}{\kappa\sqrt{t}}+\frac{2\pi\beta\delta_0^2}{\kappa^2}<\frac{\delta_0}{\kappa\sqrt{t}}+\frac{\delta_0}{\kappa\sqrt{\bar t}},
\label{align:ch15-15314}
\end{align}
与 $t\leqslant\bar t$ 矛盾。

估计\eqref{align:ch15-1532}与\eqref{align:ch15-1533}由类似论证建立。读者可在\cref{section:ch15-注记}所引文献中找到细节。

由\eqref{align:ch15-1531}、\eqref{align:ch15-1532}与\eqref{align:ch15-1533}，\eqref{align:ch15-1524}在 $t=\bar t=(\alpha\kappa\delta_0)^{-2}$ 处成立。当\eqref{align:ch15-1523}在更长的区间 $(0,T)$ 上成立、$\bar t<T\leqslant\infty$ 时，\eqref{align:ch15-1524}对任意 $t\in[\bar t,T)$ 成立，因为时间原点可平移到点 $t-\bar t$。

最后，设\eqref{align:ch15-1534}成立，并设\eqref{align:ch15-1531}在某时间区间 $(0,\hat t]$ 上成立。则对任意 $t\in(0,\hat t]$，\eqref{align:ch15-1539}、\eqref{align:ch15-1538}与\eqref{align:ch15-15311}联合蕴含
\begin{align}
\|\partial_xU(\cdot,t)\|_{L^1(-\infty,\infty)}\leqslant\frac{\delta_0}{2}+\frac{2\beta\delta_0}{\kappa^2}\int_0^t\frac1{\sqrt{\tau}}\|\partial_xU(\cdot,\tau)\|_{L^1(-\infty,\infty)}\,\mathrm{d}\tau.
\label{align:ch15-15315}
\end{align}
由 Gronwall 引理，
\begin{align}
\|\partial_xU(\cdot,t)\|_{L^1(-\infty,\infty)}\leqslant\frac{\delta_0}{2}\exp\left[\frac{4\beta\delta_0\sqrt{t}}{\kappa^2}\right],\qquad t\in[0,\hat t].
\label{align:ch15-15316}
\end{align}
因此，当 $\alpha$ 充分小但与 $\delta_0$ 无关时，\eqref{align:ch15-1531}与\eqref{align:ch15-1523}在区间 $(0,\bar t]$ 上同时成立。证明完成。
\end{proof}

\end{document}
